\section{Algebraic-rank moments}\label{sec:algebraic}

We now prove Corollary~\ref{cor:algebraic-moments}.
The inputs are the companion density and rank-equality theorem and
the exponential-moment estimate of Koymans and Smith.
The latter bounds the contribution of the density-zero exceptional set
for every fixed exponential parameter or moment order; the analytic
tail theorem is not needed.

\begin{proof}[Proof of Corollary~\ref{cor:algebraic-moments}]
Put \(N(Y)=\#\mathcal D(Y)\), \(r_d=r(E^{(d)})\), and let
\(N_j^{\mathrm{alg}}(Y)\) count the parameters with \(r_d=j\).
The rank-equality conclusion of the companion density theorem
\cite[Theorem~1.2]{OpenAIAnalyticDensityV2}, together with its analytic
rank-zero and rank-one densities, gives
\[
 \frac{N_j^{\mathrm{alg}}(Y)}{N(Y)}\longrightarrow\frac12
 \qquad(j=0,1).
\]
Indeed, analytic and algebraic ranks agree outside a density-zero set,
so their rank-\(j\) counts differ by \(o(N(Y))\). Consequently
\[
 \mathcal B(Y)=\{d\in\mathcal D(Y):r_d\ge2\}
 \quad\text{satisfies}\quad
 \#\mathcal B(Y)=o(N(Y)).
\]

For the next bound, \(E^{(v)}\) for any nonzero integer \(v\) denotes
the twist determined by its squareclass, so the integer sum may repeat
squareclasses. Theorem~1.4 of Koymans and Smith
\cite{KoymansSmith2026}, specialized to \(F=\Q\), \(A=E\),
one variable, and \(P(u)=u\), gives a constant \(C_E>0\) such that
for every fixed \(\kappa>0\),
\[
 \frac1{2Y}\sum_{\substack{v\in\Z\\0<|v|\le Y}}
 \exp\!\bigl(\kappa\,\operatorname{rank}_{\Z}E^{(v)}(\Q)\bigr)
 \le \exp\!\bigl(\exp(C_E\kappa)\bigr)
\]
for all sufficiently large \(Y\), with the threshold allowed to depend
on \(\kappa\). Their integer box already includes both signs, and
\(P(v)\ne0\) excludes exactly \(v=0\). Restricting this nonnegative
sum to \(\mathcal D(Y)\), and using
\(N(Y)\sim2Y/\zeta(2)\), gives for each fixed \(\kappa>0\)
\begin{equation}\label{eq:algebraic-exp-bound}
 \limsup_{Y\to\infty}\frac1{N(Y)}
 \sum_{d\in\mathcal D(Y)}e^{\kappa r_d}
 \le \zeta(2)\exp\!\bigl(\exp(C_E\kappa)\bigr)<\infty.
\end{equation}

For fixed \(t>0\), Cauchy--Schwarz and
\eqref{eq:algebraic-exp-bound} at \(\kappa=2t\) give
\[
 \frac1{N(Y)}\sum_{d\in\mathcal B(Y)}e^{t r_d}
 \le
 \left(\frac{\#\mathcal B(Y)}{N(Y)}\right)^{1/2}
 \left(\frac1{N(Y)}
       \sum_{d\in\mathcal D(Y)}e^{2t r_d}\right)^{1/2}
 \longrightarrow0.
\]
For fixed \(t<0\), the same exceptional average is at most
\(\#\mathcal B(Y)/N(Y)\), since \(r_d\ge0\); for \(t=0\) the full
average is identically one. The contributions from ranks zero and one
are \(N_0^{\mathrm{alg}}(Y)/N(Y)\) and
\(e^tN_1^{\mathrm{alg}}(Y)/N(Y)\), proving the exponential limit.

Finally, for each fixed positive integer \(m\), one has
\(x^{2m}\ll_m e^x\) for \(x\ge0\). A second Cauchy--Schwarz application
and \eqref{eq:algebraic-exp-bound} at \(\kappa=1\) therefore give
\[
 \frac1{N(Y)}\sum_{d\in\mathcal B(Y)}r_d^m
 \le
 \left(\frac{\#\mathcal B(Y)}{N(Y)}\right)^{1/2}
 \left(\frac1{N(Y)}\sum_{d\in\mathcal D(Y)}r_d^{2m}\right)^{1/2}
 \longrightarrow0.
\]
The rank-zero contribution is zero and the rank-one contribution is
\(N_1^{\mathrm{alg}}(Y)/N(Y)\), proving the power-moment limit without
differentiating a limiting exponential-moment formula.
\end{proof}
